\section{Objetivo General y Objetivos Específicos}

\subsection{Objetivo General}

Desarrollar y validar un detector de la fase de acumulación en esquemas de
\emph{pump \& dump} en criptoactivos, modelando el mercado como una red de
correlaciones filtrada y aplicando teoría espectral de grafos ---mediante la
Laplaciana normalizada y el \emph{Inverse Participation Ratio} (IPR)--- para
identificar patrones de coordinación artificial sobre eventos reales etiquetados.

\subsection{Objetivos Específicos}

\begin{enumerate}

  \item Construir la red de correlaciones dinámica $G_t = (V, E, w)$,
        donde cada nodo es un token, cada arista tiene peso derivado de
        la distancia de Mantegna~\cite{mantegna1999} sobre retornos
        logarítmicos, y la red evoluciona sobre ventanas deslizantes.
        Los datos de entrada son velas de 1~minuto recopiladas mediante
        la API pública de Binance (\texttt{GET /api/v3/klines}) para los
        ${\approx}351$ eventos del dataset de La~Morgia
        \textit{et al.}~\cite{lamorgia2020}, pares \texttt{SYM/BTC},
        período 2021--2022. Binance ofrece mayor cobertura histórica que
        datasets anteriores, reduciendo el problema de tokens que ya no
        existen.

  \item Filtrar el ruido estadístico de $C_t$ mediante RMT (límites de
        Marchenko-Pastur, ventana dual $T_{\text{cal}} = 720~\text{h}$ /
        $T_{\text{det}} = 48~\text{h}$, reconstrucción $\tilde{C}$ según
        la ecuación~(\ref{eq:Ctilde})) y esparsificar la red resultante
        usando el \emph{Árbol de Expansión Mínima} (MST) como esqueleto
        topológico, al que se añaden $k$ aristas adicionales calibradas
        por modularidad para preservar la estructura de comunidades del
        grafo completo.

  \item Aplicar teoría espectral de grafos sobre $G_t$: calcular la
        Laplaciana normalizada $\mathcal{L}_t$ del grafo esparsificado y
        obtener el perfil completo del \emph{Inverse Participation Ratio}
        para todos los modos $k \geq 1$. No restringimos el análisis a
        ninguna región del espectro: la posición espectral de una anomalía
        transitoria no es predecible de antemano, a diferencia de los
        enfoques estáticos de detección de comunidades. Los estadísticos
        distribucionales que usamos son:
        \begin{align*}
          \mathrm{IPR}_{\text{mean}}(t) &= \frac{1}{N-1}
            \sum_{k=1}^{N-1} \mathrm{IPR}_k(t), \\[3pt]
          \mathrm{IPR}_{p90}(t) &= \mathrm{percentil}_{90}
            \bigl\{\mathrm{IPR}_k(t) : k = 1, \ldots, N-1\bigr\}.
        \end{align*}
        Al no depender del orden de los eigenvalores, ambos estadísticos son
        inmunes al problema del \emph{eigenvalue crossing}.

  \item Definir el umbral de detección usando el \emph{Configuration
        Model}~\cite{newman2003} como punto de comparación. Usar simplemente
        un percentil histórico como umbral es problemático: un mercado agitado
        por una crisis externa (colapso de FTX en 2022, por ejemplo) puede
        elevar el IPR sin que haya manipulación. En cambio, para cada ventana
        $t$ generamos $M = 10\,000$ redes aleatorias con la misma estructura de
        conexiones que $G_t$ (misma secuencia de grados, aristas reasignadas al
        azar) y estimamos cuánto IPR produce ese ruido de fondo. Declaramos
        una anomalía cuando $\mathrm{IPR}_{p90}(t)$ supera el percentil~99.7
        empírico de la distribución nula:
        \begin{equation}
          \mathrm{IPR}_{p90}(t) > \hat{F}_{\text{null}}^{-1}(0.997),
          \label{eq:threshold}
        \end{equation}
        donde $\hat{F}_{\text{null}}^{-1}$ es la función de cuantiles empírica
        de los $M = 10\,000$ valores simulados. Este umbral no paramétrico evita
        asumir normalidad en la distribución nula: el estadístico
        $\mathrm{IPR}_{p90}$ es un estadístico de orden sobre una métrica de
        red y su distribución bajo el Configuration Model puede ser asimétrica.
        La equivalencia con $\mu_{\text{null}} + 3\,\sigma_{\text{null}}$
        se cumple solo bajo normalidad y se reportará como referencia, pero
        el criterio operativo es el percentil empírico. Esto controla por la
        densidad local de la red: un IPR alto no puede atribuirse simplemente
        a que la red sea densa, porque eso ya está capturado por el modelo
        nulo. Reportaremos resultados bajo el modelo estándar y bajo una
        versión con pesos permutados, para evaluar robustez.

  \item Evaluar el detector sobre los ${\approx}351$ eventos confirmados del
        dataset~\cite{lamorgia2020} (pares \texttt{SYM/BTC} en Binance,
        2021--2022). Por el severo desbalance de clases propio de la detección
        de anomalías, la métrica principal es el \textbf{PR-AUC} (área bajo la
        curva Precisión-Recall), complementada con el \textbf{F1-Score}. El
        ROC-AUC se reporta para comparabilidad con la literatura, pero no es la
        métrica de decisión porque los Verdaderos Negativos inflan su
        denominador de forma engañosa en datasets desbalanceados. Además,
        inspeccionaremos los eigenvectores localizados para verificar que los
        tokens con mayor peso en $\boldsymbol{\phi}_{k^*}(t)$ corresponden
        efectivamente a los tokens manipulados, dando interpretabilidad al
        resultado.

\end{enumerate}
